\documentclass[11pt,a4paper]{article}
\usepackage[margin=2.5cm]{geometry}
\usepackage{amsmath,amssymb,booktabs,xcolor,enumitem}
\usepackage[normalem]{ulem}
\usepackage[hidelinks]{hyperref}
\setlength{\parskip}{0.6em}\setlength{\parindent}{0pt}
\definecolor{newc}{HTML}{1f6f43}
\newcommand{\new}[1]{\textcolor{newc}{#1}}
\newcommand{\old}[1]{\textcolor{gray}{\sout{#1}}}
\newcommand{\oldmath}[1]{\textcolor{gray}{#1}}
\newcommand{\mockup}[1]{\textit{#1}}
\title{Proposed change [A93], congestion is measured as the time picks stand still at each shelf}
\date{13 November 2026, version 3 (plainer words, the reason for keeping demand on the nearest shelf)}
\begin{document}
\maketitle

\section*{How to read this}
Part~1 says what we found and what is proposed. Part~2 shows the thesis text that
would change. \new{Green} is new text, \old{struck} text goes, black text stays,
and italic text is a note on what to draw or add. A removed equation is grey and
marked ``removed''. Part~3 has the questions. Nothing is changed in the thesis or
the code yet.

\section*{Part 1. What we found}

\subsection*{The old signal did not work}
Congestion-aware storage counts, for each edge, how often agents passed it times
how often they waited to get onto it.

\begin{itemize}[leftmargin=1.2em]
\item \textbf{The number is far too small.} It is about $10^{-4}$ per edge, while a
  trip is about 29 steps. So it never changed which shelf was chosen. This is why
  $\beta$ from 0 to 4 made almost no difference.
\item \textbf{It jumps around.} Only 18 to 34 of 782 edges have any wait in a
  storage period. The order of the shelves it gives changes almost completely from
  one period to the next.
\item \textbf{It looks in the wrong place.} With Token Passing a pick stands still
  only 1 to 2.5 steps. Most of the time is spent waiting for a free agent, and that
  wait is not on any edge.
\end{itemize}

We tried several other edge signals with Token Passing on 10 seeds. None of them
was better than demand storage.

\subsection*{The learned controllers do get stuck}
The learned agents stand still much longer than with Token Passing, and more the
bigger the fleet.

\begin{center}
\begin{tabular}{lcc}
\toprule
 & steps standing still per pick & Token Passing \\
\midrule
8 agents, decentralised & about 21 & 1 to 2.5 \\
16 agents, section & about 32 & 1 to 2.5 \\
16 agents, decentralised & about 38 & 1 to 2.5 \\
32 agents, decentralised & 74 to 85 & 1 to 2.5 \\
\bottomrule
\end{tabular}
\end{center}

The agents stand still in queues, when two agents face each other, and at the
workstation lock of [A91]. Counting waits on edges misses most of this.

\subsection*{The new signal}
When a pick is done, we count the steps its agent stood still and note them on the
shelf it came from. The signal of a shelf is the average over its recent picks.

\begin{itemize}[leftmargin=1.2em]
\item It is counted in steps, like the trip. So $\beta = 1$ means one step standing
  still counts the same as one step of travel.
\item It catches every reason for standing still.
\item It belongs to a shelf, which is exactly what storage chooses.
\end{itemize}

At 8 agents, shelves differ a lot, from about 3 to 68 steps per pick, and the order
of the shelves stays almost the same from one period to the next.

The signal is used in two places. When stock arrives, it goes where picks rarely
stand still. And when a product is on several shelves, a pick takes it from the
shelf with the lowest distance plus standing still, not just the nearest one. The
second happens for every order, so the loop now also reacts quickly.

\subsection*{The test}
We used policies that are already trained under demand storage, the best saved
policy of each run. Nothing was retrained. We compared demand storage as it is with
the signal used when stock arrives, when picking, and in both. 5 seeds, 4000 steps,
of which the first 1000 are not counted. The test also averaged the demand numbers
over periods (question~2).

\begin{center}
\begin{tabular}{llccccc}
\toprule
 & & steps still & less still & tasks & as many & keeps \\
 & & per pick & than today & done & tasks or more & up \\
\midrule
8 agents & today & 21.0 & & 339 & & 4/5 \\
decentralised & when stock arrives & 18.6 & 3/5 & 339 & 3/5 & 4/5 \\
 & when picking & 20.8 & 3/5 & 334 & 3/5 & 4/5 \\
 & \textbf{both} & \textbf{17.3} & \textbf{5/5} & 340 & 3/5 & 4/5 \\
\midrule
16 agents & today & 32.5 & & 361 & & 0/5 \\
section & when stock arrives & 29.4 & 4/5 & 385 & 4/5 & 0/5 \\
 & when picking & 31.6 & 4/5 & 374 & 3/5 & 0/5 \\
 & \textbf{both} & \textbf{29.1} & \textbf{5/5} & \textbf{384} & 4/5 & 0/5 \\
\bottomrule
\end{tabular}
\end{center}

Used in both places, the signal cut the time standing still on every seed. It was
18 percent less at 8 agents and 10 percent less at 16. At 16 agents the agents also
did 6 percent more tasks. At 8 agents a pick took 244 steps instead of 294. This is
the first congestion signal that helps.

Two runs tell us nothing. The 16-agent decentralised controller does only 115 to
150 tasks, where the section controller does 360 to 385. The 32-agent controller
does 4 to 9 tasks in 3000 steps. Both get stuck over a long run. You said the
16-agent decentralised controller needs messages to work, so we test it with
messages (question~5).

\subsection*{What this test is not}
It is a short test on 5 seeds, with policies that never trained with the signal. It
shows the signal is worth keeping. The full test in question~5 gives the result.
Sending orders to the least busy workstation and moving old stock were tried the
same day. They are not part of this change (question~7).

\section*{Part 2. The text}

\subsection*{2.1 Problem Formulation, ``What the rule knows''}
\begin{quote}
\old{The rule only knows the last storage period. It reads three averages per
step over that period:}
\new{The rule only knows the recent past. It reads two numbers:}
\begin{itemize}
\item \(\hat\rho_t(k)\), how often SKU \(k\) was ordered \new{per step in the last storage period}
\item \old{\(\hat\mu_t(e)\), how often agents crossed edge \(e\)}
\item \old{\(\hat w_t(e)\), how many steps agents waited to get onto edge \(e\)}
\item \new{\(\hat\omega_t(v)\), how many steps recent picks from shelf \(v\) stood still, on average}
\end{itemize}
Demand storage uses only the first, the orders. Congestion-aware storage uses
\old{all three. The last two come from how the agents moved, and show which
aisles were crowded. An edge counts as crowded when many agents cross it and
they often wait to get onto it.}
\new{both. The second shows where agents lose time because of each other. It may
be a queue in an aisle, two agents facing each other, or a queue at a
workstation.}
\end{quote}

In \texttt{eq:storageupdate} the inputs change from
\oldmath{$\hat\rho_t,\hat\mu_t,\hat w_t$ (removed)} to \new{$\hat\rho_t,\hat\omega_t$}.

\subsection*{2.2 Problem Formulation, ``What the rule decides''}
\begin{quote}
Putting every popular SKU on the best shelves would send all traffic down
the same few aisles. Congestion-aware storage therefore also avoids shelves
\old{whose trip runs through crowded aisles.}
\new{where picks often stand still.}

\ldots

The crowding cost is this thesis's own, because the classic problem has no
traffic. It adds up, over the same units, \old{the crowding $\hat\mu_t(e)\,\hat w_t(e)$
on every edge $e$ of the trip through $v$,}
\new{how often their SKU is ordered times how long picks stand still at their
shelf. Every pick from that shelf meets the wait, so a popular SKU pays it more
often.}
\begin{equation*}
\oldmath{K = \sum_{k}\sum_{v}\Bigl(\sum_{e \in \mathrm{path}(v)} \hat\mu_t(e)\,\hat w_t(e)\Bigr)\, x^{\mathrm{new}}_t(k,v) \qquad \text{removed}}
\end{equation*}
\begin{equation*}
\color{newc} K = \sum_{k\in\mathcal{K}} \sum_{v\in V_{\mathrm{str}}}
\hat\rho_t(k)\,\hat\omega_t(v)\, x^{\mathrm{new}}_t(k,v) .
\end{equation*}
The rule aims for a placement with a small value of $D + \beta K$,
where $\beta$ sets how much \old{crowding} \new{standing still} counts against distance.
\new{Both count steps, so with $\beta = 1$ a step standing still counts as a step of
travel.}
\end{quote}

\subsection*{2.3 Problem Formulation, the pickup rule}
\begin{quote}
A line takes its units from the shelf nearest its workstation first. If that shelf
holds enough, the line is one task. \ldots\ This is the \emph{pickup rule}.

\new{Under congestion-aware storage a line chooses differently. For each shelf it
adds the distance to the workstation and the time picks there stand still,
$\hat\omega_t(v)$. It takes from the shelf with the smallest sum first. Say tea is
on $A$ and $C$. Picks from $A$ have stood still 5 steps on average and picks from
$C$ 1 step. Then the line goes to $C$ if $C$ is at most 3 steps farther away.}

\new{Fixed and demand storage keep taking the nearest shelf. Demand storage knows
only the orders, so letting it pick by standing still would make it read the
congestion too. The gap between demand and congestion-aware storage would then no
longer show what reading congestion adds. The nearest shelf first is also the plain
rule a warehouse uses when it knows nothing about traffic. So only the rule that
closes the loop changes how lines pick.}

\new{This choice is made for every order line, while stock only arrives every
$\Delta$ steps. So congestion-aware storage reacts both slowly and quickly.}
\end{quote}

\mockup{Appendix~C, one line after eq:pickuprule: the shelves $v_1, v_2, \dots$
that hold free units of $k$ are sorted by $d_G(v_i, g)$, or by
$d_G(v_i, g) + \hat\omega_t(v_i)$ under congestion-aware storage, and by vertex
number when equal.}

\subsection*{2.4 Problem Formulation, the loop}
\begin{quote}
Backward, routing shapes storage. \old{The traffic the agents create is measured
in the estimates $\hat\mu_t$ and $\hat w_t$, and at the next epoch the storage
update uses them to decide where items go.}
\new{How long agents stand still is measured at each shelf in $\hat\omega_t$. At the
next storage update it decides where new stock goes. Every new order line also
uses it to choose the shelf it picks from.}
With fixed storage only the forward path exists. The backward path is
what congestion-aware storage adds, because only that rule reads the
\old{traffic.} \new{time agents stand still.}
\end{quote}

\mockup{Figure loop.pdf, panel (b): the backward arrow says $\hat\omega_t$ instead
of $\hat\mu_t, \hat w_t$, and a short second arrow goes from $\hat\omega_t$ to the
pickup of each task. Caption: ``The two-timescale loop under fixed (a) and
congestion-aware storage (b). Demand storage reads only $\hat\rho_t$, so standing
still never reaches it.''}

\subsection*{2.5 Implementation, ``Where arriving stock goes''}
\begin{quote}
\textbf{Congestion-aware storage} ($F_{\mathrm{cng}}$) adds \old{crowding}
\new{standing still} to the trip. Take two shelves with the same trip distance,
\old{one reached through a quiet aisle and one through an aisle where agents often
queue.} \new{one where picks rarely stand still and one where they often do.}
Demand storage treats them as equal, and congestion-aware storage prefers the
quiet one. Formally, the \emph{congestion-aware score} of SKU $k$ at shelf $v$ is
\begin{equation*}
\oldmath{\mathrm{score}(k,v) = \hat\rho_t(k)\cdot d^{\mathrm{trip}}_t(k,v) + \beta \sum_{e \in \mathrm{path}(v)} \hat\mu_t(e)\,\hat w_t(e) \qquad \text{removed}}
\end{equation*}
\begin{equation*}
\color{newc} \mathrm{score}(k,v) = \hat\rho_t(k)\,\bigl(d^{\mathrm{trip}}_t(k,v) + \beta\,\hat\omega_t(v)\bigr) ,
\end{equation*}
\old{where $\mathrm{path}(v)$ is the trip through $v$.} The weight $\beta \ge 0$
sets how much \old{crowding} \new{standing still} counts against distance.

\old{The estimates are the per step rates of eq:estimates, taken over the last
storage period only.}
\new{The demand number is the rate per step of eq:estimates over the last storage
period.} Counting per step keeps $\beta$ comparable across the sweep of the
storage period $\Delta$. If no order was released in a period, nothing was picked
and nothing arrives.

\new{The time standing still comes from the picks done in the period. For each, we
count the steps its agent stood still while serving it. This is the blocked part
of its service time in eq:split. The average over a shelf's picks is the new
reading for that shelf. A shelf gets only a few picks in one period, about eight at
8 agents, so one reading alone jumps around. So the number keeps four fifths of its
old value and adds one fifth of the new reading. A shelf with no pick in the period
keeps its old value. A shelf's first reading is used as it is. With this, the order
of the shelves hardly changes from one period to the next.}

\old{A wait is counted on the edge the agent is waiting to enter, the next edge of
its shortest route to its goal. It counts whenever an agent with a task stands
still, whether it chose to wait or conflict resolution stopped it. Counting only
the stops would show no waiting for Token Passing, which never collides but waits
on purpose to let others pass. An agent without a task, or already at its goal,
adds no wait.}
\end{quote}

\subsection*{2.6 Appendix C, eq:estimates}
\begin{equation*}
\hat\rho_t(k) = \tfrac{1}{\Delta}\,\#\{\text{order lines for } k \text{ released in the window}\},
\end{equation*}
\begin{equation*}
\oldmath{\hat\mu_t(e) = \ldots,\quad \hat w_t(e) = \ldots \qquad \text{removed}}
\end{equation*}
\begin{equation*}
\color{newc} \hat\omega_t(v) = \tfrac{4}{5}\,\hat\omega_{t-\Delta}(v) + \tfrac{1}{5}\,
\frac{1}{|P_t(v)|}\sum_{j \in P_t(v)}\;\sum_{t'=t_j^{\mathrm{assigned}}}^{t_j^{\mathrm{finish}}-1}\omega_j(t') ,
\end{equation*}
\new{where $P_t(v)$ are the picks from shelf $v$ done in the window. If there are
none, $\hat\omega_t(v) = \hat\omega_{t-\Delta}(v)$. A shelf's first reading is used
as it is.}

\subsection*{2.7 Appendix G, ``Calibration: the congestion weight''}
\begin{quote}
The congestion weight is $\beta = 1$. Both terms of the score
count steps, \old{the expected travel and the expected waiting, so a step of
waiting weighs the same as a step of travel.}
\new{travel and standing still, so a step standing still counts the same as a step
of travel.}

\new{The rule first counted, for each edge, passes times waits. With Token Passing
a weight from 0 to 4 changed the pick service time by at most 2 percent. The number
was only about $10^{-4}$ per edge, and a pick stood still only 1 to 2.5 steps.
Other edge signals did no better on 10 seeds. The time picks stand still at each
shelf took its place [A93].}

\new{The learned controllers stand still much longer. With them the new signal cut
the time standing still on every seed at 8 agents, and at 16 agents with the
section controller, by 18 and 10 percent.} \mockup{Table from Part~1, with the
final numbers from the full test of question~5.}

\old{The first sweep, on the first version of the warehouse graph, \ldots\ that
wait lies on no edge the weight can see.} \mockup{The three old paragraphs go, and
their figure app\_storage\_beta.png too.}
\end{quote}

\subsection*{2.8 Notation table and parameters}
\begin{itemize}[leftmargin=1.2em]
\item \old{\(\hat\mu_t(e)\), traffic estimate} and \old{\(\hat w_t(e)\), waiting estimate} go.
\item \new{\(\hat\omega_t(v)\), time standing still at a shelf. The average steps
  recent picks from shelf \(v\) stood still, kept over storage periods. Used by
  \(F\) and, under congestion-aware storage, by the pickup rule.}
\item \(\beta\), congestion weight. How much the \old{crowding} \new{standing still}
  cost \(K\) counts against the slotting cost \(D\).
\item Parameters table, the reason for \(\beta = 1\): \old{a step of waiting weighs
  as a step of travel} \new{a step standing still counts as a step of travel}.
  \new{New row: the newest period counts one fifth in $\hat\omega_t$.}
\end{itemize}

\section*{Part 3. Questions}
\begin{enumerate}[leftmargin=1.4em]
\item \textbf{The new pick choice only for congestion-aware storage?} Decided yes
  (13 November). The reason is in 2.3. Demand storage knows only the orders, so the
  gap to congestion-aware storage must come from reading congestion alone. The main
  result, demand against fixed, stays the same and needs no new runs.
\item \textbf{Average the demand numbers too?} The test also averaged
  $\hat\rho_t$ over periods. To match the test, congestion-aware storage should do
  the same. But then the difference to demand storage comes from two changes. With
  Token Passing the averaging alone was 1.5 percent faster (8 of 10 seeds,
  $p = 0.08$). The other choice is to test again without it.
\item \textbf{One fifth for the newest period.} It was picked, not tuned. Keep it,
  or try a few values in the full test?
\item \textbf{Remove $\hat\mu_t$ and $\hat w_t$ completely?} I suggest yes, and also
  the paragraph on which edge a wait is counted on ([A36]). Nothing else uses them.
\item \textbf{The full test.} 20 test seeds, 10000 steps with the first 5000 not
  counted, the best saved policy of each run, demand against congestion-aware
  storage. Runs: 8 agents decentralised, 16 agents section, and 16 agents
  decentralised with messages. The 32-agent runs wait for a policy that does not
  get stuck. Retrain under the new rule later, or not at all?
\item \textbf{The name.} The rule now also chooses where to pick. Keep
  ``congestion-aware storage'', or say once that it also picks?
\item \textbf{Least busy workstation and moving old stock.} The least busy
  workstation made Token Passing about 28 percent faster on two seeds but did
  nothing for the learned controller. Moving old stock made Token Passing slower.
  Leave both out of [A93]?
\end{enumerate}

\end{document}
